\chapter{Introdução}

\section{Considerações Iniciais}

A área de desenvolvimento de software tem evoluído com a necessidade de sistemas mais
escaláveis e de fácil manutenção, tornando a arquitetura de software um fator essencial no
desenvolvimento de aplicações. As decisões arquiteturais definem a organização dos elementos
de um sistema e influenciam diretamente atributos de qualidade como desempenho,
escalabilidade e manutenibilidade \cite{bass2021}. Entre os principais modelos arquiteturais
adotados no desenvolvimento de aplicações Web, destacam-se o monolítico e o de microsserviços
\cite{softdesign2020}.

A arquitetura monolítica concentra toda a aplicação em uma única unidade de implantação, o
que tende a facilitar o desenvolvimento, os testes e a implantação nas fases iniciais de um
projeto, mas pode dificultar a escalabilidade e a manutenção à medida que o sistema cresce
\cite{richardson2018}. Já a arquitetura de microsserviços divide o sistema em serviços
pequenos e independentes, cada um executado em seu próprio processo e implantado de forma
autônoma, comunicando-se por mecanismos leves, como requisições HTTP ou mensageria
\cite{lewis2014}. Essa abordagem oferece maior flexibilidade para evoluir e escalar partes
específicas do sistema, porém transfere para a aplicação a complexidade inerente aos sistemas
distribuídos, como latência de rede, falhas parciais e consistência de dados entre serviços
\cite{newman2021}.

Nos últimos anos, os microsserviços passaram a ser amplamente difundidos na indústria, muitas
vezes apresentados como a escolha natural para sistemas modernos. Entretanto, a adoção desse
estilo nem sempre traz os benefícios esperados: há relatos de organizações que, diante de
custos operacionais elevados, reagruparam serviços distribuídos em uma única aplicação, como
no caso do serviço de monitoramento de vídeo do Prime Video \cite{kolny2023}. Autores como
\citeonline{fowler2015} e \citeonline{newman2021} recomendam, inclusive, que sistemas novos
comecem como monólitos e só sejam decompostos quando houver razões concretas para isso.

Diante desse cenário, surge a seguinte questão: em quais aspectos e em quais contextos cada
uma dessas arquiteturas se mostra mais vantajosa? Para respondê-la, este trabalho compara as
duas abordagens por meio do desenvolvimento de duas aplicações de \textit{e-commerce} com as
mesmas tecnologias e funcionalidades, uma em arquitetura monolítica e outra em arquitetura de
microsserviços, ambas implementadas em Java com o \textit{framework} Spring Boot. As duas
versões oferecem cadastro e autenticação de usuários, catálogo de produtos, personalização da
loja, realização de pedidos e processamento de pagamentos integrado ao Mercado Pago, e são avaliadas segundo os critérios de
desempenho, escalabilidade, manutenibilidade, complexidade de desenvolvimento e implantação.
Como resultado, espera-se identificar as principais vantagens e desvantagens de cada
abordagem, contribuindo para a escolha da arquitetura mais adequada em diferentes contextos.

\section{Objetivos do Trabalho}

Como objetivo geral, serão desenvolvidas duas aplicações de \textit{e-commerce} com as mesmas
funcionalidades, utilizando as arquiteturas monolítica e de microsserviços, a fim de comparar
suas características, vantagens e limitações. Para alcançar esse objetivo, foram definidos os
seguintes objetivos específicos:

\begin{itemize}
    \item Analisar as diferenças conceituais entre a arquitetura monolítica e a arquitetura de
    microsserviços;
    \item Levantar o estado da arte sobre a comparação entre as duas arquiteturas, identificando
    os métodos, as métricas e os resultados de trabalhos anteriores;
    \item Desenvolver aplicações com as mesmas funcionalidades em ambas as arquiteturas;
    \item Avaliar as duas aplicações quanto ao desempenho, à escalabilidade, à
    manutenibilidade, à complexidade de desenvolvimento e à implantação;
    \item Identificar os benefícios e desafios de cada abordagem;
    \item Auxiliar desenvolvedores na escolha da arquitetura mais adequada para cada projeto.
\end{itemize}

\section{Justificativa}

Com o avanço das aplicações modernas, a escolha da arquitetura de software tornou-se essencial
para garantir desempenho, escalabilidade e manutenção adequados. Por envolver decisões
difíceis de reverter depois que o sistema está em produção, a arquitetura exige que suas
escolhas sejam fundamentadas na análise de \textit{trade-offs}, isto é, das vantagens e
desvantagens de cada alternativa \cite{richards2024}. Nesse contexto, compreender diferentes
modelos arquiteturais é fundamental para a formação de profissionais da área.

Atualmente, observa-se um aumento na adoção de microsserviços, embora nem sempre essa escolha
seja a mais adequada, devido à complexidade que acrescenta ao desenvolvimento e à operação dos
sistemas \cite{newman2021, aws-monolito-microsservicos}. Os estudos que comparam
experimentalmente as duas arquiteturas apresentam resultados divergentes, que variam conforme
a aplicação, a carga e o ambiente de execução \cite{aldebagy2018, blinowski2022}. Diante disso,
surge a necessidade de analisar, de forma prática e controlada, quando utilizar microsserviços
ou monolito.

Assim, este trabalho se justifica pela proposta de comparar, na prática, essas duas
arquiteturas em uma aplicação com funcionalidades realistas, como autenticação, pagamento
integrado, armazenamento de imagens e comunicação assíncrona, contribuindo para uma melhor
tomada de decisão no desenvolvimento de sistemas.

\section{Estrutura da Monografia}

Este trabalho está organizado em sete capítulos. O Capítulo~\ref{cap:referencial} apresenta o
referencial teórico sobre arquitetura de software, as arquiteturas monolítica e de
microsserviços, as formas de comunicação entre componentes e os atributos de qualidade
utilizados na comparação. O Capítulo~\ref{cap:estado-arte} apresenta o estado da arte, descrevendo
os trabalhos que compararam as duas arquiteturas e posicionando este trabalho em relação a
eles. O Capítulo~\ref{cap:metodologia}
descreve a metodologia adotada, incluindo os requisitos das aplicações, as tecnologias, o
ambiente de testes e as métricas de comparação. O Capítulo~\ref{cap:desenvolvimento} detalha o
desenvolvimento das duas versões da aplicação de \textit{e-commerce} e as diferenças entre
elas. O Capítulo~\ref{cap:resultados} apresenta e discute os resultados obtidos em cada
critério de comparação. Por fim, o Capítulo~\ref{cap:conclusoes} traz as conclusões, as
contribuições do trabalho e as sugestões de trabalhos futuros.
